% Draft section for WILSON_FIRST_RG_STEP_V4.tex: two disjoint squares and the cross terms.
% Notation and labels are those of WILSON_FIRST_RG_STEP_V3.tex (hash A520A59F363D).
% Fragment, not compiled alone.

\section{Two Disjoint Squares}\label{sec:two-loops}

This section extends the extraction of Section~\ref{sec:joint-extraction} and Appendices~\ref{app:signed-defects} and~\ref{app:regional-scalar} from one square to the product of the Wilson functions of two disjoint squares in parallel coordinate planes. Nothing in the single-square sections is changed; the objects of this section are defined for the pair and compared with the single-square objects of each square.

\subsection{Setting}\label{sec:two-loops-setting}
Let $C_1,C_2$ be coarse translates of the counterclockwise squares with vertices $0,n_iLe_1,n_iL(e_1+e_2),n_iLe_2$, $i=1,2$, in fine-lattice units (so both lie in planes parallel to that of $e_1,e_2$); let $V(C_i)$ be the set of fine sites of $C_i$, the endpoints of its bonds, and assume $V(C_1)\cap V(C_2)=\varnothing$. Let $\chi_1,\chi_2$ be normalised characters of irreducible unitary representations $\pi_1,\pi_2$ of $G$, $\chi_i=\dim(\pi_i)^{-1}\operatorname{tr}\pi_i$. The reversed orientation of $C_i$ is the case $\bar\chi_i$, since $\chi_i(U^{-1})=\overline{\chi_i(U)}$. Put $C=C_1\sqcup C_2$ and
\begin{equation}\label{eq:two-loop-observable}
    W_{\chi_1,\chi_2}(C;U)=\chi_1\bigl(\Hol_UC_1\bigr)\,\chi_2\bigl(\Hol_UC_2\bigr)\;.
\end{equation}
The traversal set is $\mathcal I_C=\mathcal I_{C_1}\sqcup\mathcal I_{C_2}$, the traversals of each $C_i$ numbered cyclically from a base point of $C_i$; $n_{C_i}=4n_iL$ and $n_C=n_{C_1}+n_{C_2}$. Each square satisfies the square/torus geometry of Section~\ref{sec:extent-weights} with its own side, $K_\mu\geq3(n_i+2)$ for every coarse period and $i=1,2$; no condition relating the positions of the two squares is imposed beyond $V(C_1)\cap V(C_2)=\varnothing$. For $a\in\mathcal I_{C_i}$ let $\delta_a$ be the periodic $\ell^\infty$ fine distance from the initial site of the traversal $a$ to $V(C_{3-i})$, and put
\begin{equation}\label{eq:two-loop-cross-sum}
    \Delta_\times=\sum_{a\in\mathcal I_C}e^{-c_\times\delta_a/(LM)}\leq n_Ce^{-c_\times d/(LM)}\;,\qquad d=\min_{a\in\mathcal I_C}\delta_a\;,
\end{equation}
with the constant $c_\times>0$ of Lemma~\ref{lem:two-loop-cross}, fixed before the coupling.

A superscript or subscript $i$ marks the single-square objects of the square $C_i$ with the character $\chi_i$: the contact sets $\mathsf T^i_D=\mathsf T_D\cap\mathcal I_{C_i}$, the groups of Definition~\ref{def:groups} for the contour $C_i$ alone, their reference scalars $a^{(i)}_W$ of \eqref{eq:part1-reference-group}, the reference gas $S^\circ_{\xi,C_i}$ of \eqref{eq:part1-reference-gas}, and the single-square lemmas of Appendix~\ref{app:signed-defects} with $C=C_i$. All of them are formed on the same branch $\xi$ with the same carrier, cubes, hosts, measures and cluster occurrences: the domains, probe measures $\widehat\nu^{\mathrm f}_{\xi,D}$, vacuum clusters and reference data of Sections~\ref{sec:fixed-first-kernel}--\ref{sec:joint-extraction} do not involve the observable, which enters only through the contact sets, $f_\xi$, $f^\circ_\xi$ and the reference scalars $a_W$. The coupling is below the thresholds of the single-square results for $\chi_1$ and for $\chi_2$ and below those of this section.

Covered and far are taken for both squares: $\xi\in A_{\mathrm{cov}}$ means that $\xi$ is covered (Lemma~\ref{lem:covered}) for the contour $C=C_1\sqcup C_2$, which, the conditions being imposed bond by bond, is the intersection of the covered events of $C_1$ and $C_2$; a history is far when $\operatorname{dist}(\Gamma_C,E(\Lambda_1^c))>b$ with $\Gamma_C=\Gamma_{C_1}\cup\Gamma_{C_2}$, that is far for both squares.

\subsection{Groups for two squares}\label{sec:two-loops-groups}
For the pair, $f_\xi(a)=W_{\chi_1,\chi_2}(C;U_\xi(a))$, and the contact set $\mathsf T_D$ of \eqref{eq:part1-contact-set} is taken over $n\in\mathcal I_C$. The argument preceding Definition~\ref{def:groups} applies verbatim up to the formation of the components: unmarked domains contribute zero, $q_\xi(\mathcal A)$ is expanded by \eqref{eq:c-cluster-expansion}, and the graph on the marked roots and the cluster occurrences of a term has components $\gamma$, each with a marked root and the contact set $\mathsf T_\gamma\subseteq\mathcal I_C$, the union of the contact sets of its roots. For $i=1,2$ let $H^i_\gamma=\varnothing$ if $\mathsf T_\gamma\cap\mathcal I_{C_i}=\varnothing$; otherwise $H^i_\gamma$ is the hull of $\mathsf T_\gamma\cap\mathcal I_{C_i}$ in the cyclic set $\mathcal I_{C_i}$, defined as in Section~\ref{sec:joint-extraction}: the shortest cyclic interval containing it when that interval has fewer than $n_{C_i}/2$ positions, and $\mathcal I_{C_i}$ otherwise. Put $H_\gamma=H^1_\gamma\sqcup H^2_\gamma$. Two components overlap when $H^i_\gamma\cap H^i_{\gamma'}\neq\varnothing$ for some $i$; the overlap classes are the connected components of the overlap graph on the components of the term. A set $K$ of components is long on $C_i$ when $\bigl|\bigcup_{\gamma\in K}H^i_\gamma\bigr|\geq n_{C_i}/2$.

\begin{definition}[groups for two squares]\label{def:two-loop-groups}
A partition $\mathfrak P$ of the components of a term is admissible if (A1) every overlap class lies in one block, and (A2) a block that is long on $C_i$ contains every component $\gamma$ with $\mathsf T_\gamma\cap\mathcal I_{C_i}\neq\varnothing$. The groups of the term are the blocks of its finest admissible partition (Lemma~\ref{lem:two-loop-grouping}(a)). A group $W$ consists of its root family $\mathcal A(W)$, its ordered list of cluster occurrences, and its interval set $\mathcal H(W)=\mathcal H_1(W)\sqcup\mathcal H_2(W)$, where $\mathcal H_i(W)=\mathcal I_{C_i}$ if $W$ is long on $C_i$, read from the base point of $C_i$, and $\mathcal H_i(W)=\bigcup_{\gamma\in W}H^i_\gamma$ otherwise, a finite union of cyclic intervals (a cyclic interval or $\mathcal I_{C_i}$ for a $C_i$-group, by Lemma~\ref{lem:two-loop-grouping}(d)); as in Definition~\ref{def:groups}, these data arise from a term, and conversely the term with root family $\mathcal A(W)$ and cluster list of $W$ has $W$ as its only group, with the same interval set (Lemma~\ref{lem:two-loop-grouping}(a)). Its weight $w_\xi(W)$ is \eqref{eq:part1-group-weight}. $W$ is a $C_i$-group if $\mathcal H_{3-i}(W)=\varnothing$ and a cross group if $\mathcal H_1(W)\neq\varnothing\neq\mathcal H_2(W)$. Two groups are compatible when all their constituent domains are mutually compatible and $\mathcal H_i(W)\cap\mathcal H_i(W')=\varnothing$ for $i=1,2$. For a compatible group family, $\mathcal A(\mathcal S)$ and $A_\xi(\mathcal S)$ are \eqref{eq:part1-closed-family-pairing}.
\end{definition}

\begin{lemma}[two-square grouping]\label{lem:two-loop-grouping}
(a) Every term has a finest admissible partition; each of its blocks that is long on neither square is an overlap class. If no component of the term has a contact on $C_{3-i}$, it is the grouping of Section~\ref{sec:joint-extraction} for the square $C_i$. For a group $W$ of a term, the term $T_W$ with root family $\mathcal A(W)$ and the cluster list of $W$ has $W$ as its only group, with the same interval set. (b) The groups of a term are pairwise compatible, and Lemma~\ref{lem:grouped-source} holds for the pair: $F_\xi(0)=\sum_{\mathfrak B\text{ compatible}}A_\xi(\mathfrak B)$, absolutely convergent. (c) A $C_i$-group of the pair is a group of the single square $C_i$ (Definition~\ref{def:groups} with $C=C_i$) none of whose roots has a contact on $C_{3-i}$; conversely every such single-square group is a $C_i$-group of the pair, with the same interval, weight and footprint (the footprint of a group of the pair is defined in Section~\ref{sec:two-loops-pins}). Two $C_i$-groups are compatible for the pair if and only if they are compatible for the square $C_i$; a $C_1$-group and a $C_2$-group are compatible if and only if their constituent domains are mutually compatible. (d) If $W$ is a $C_i$-group that is not long on $C_i$, then $\mathcal H_i(W)$ is a cyclic interval of fewer than $n_{C_i}/2$ traversals.
\end{lemma}

\begin{proof}
(a) The one-block partition is admissible. If $\mathfrak P$ and $\mathfrak Q$ are admissible, so is their common refinement, whose blocks are the nonempty sets $P\cap Q$: (A1) holds because every overlap class lies in one block of each; for (A2), if $P\cap Q$ is long on $C_i$, then so are $P\supseteq P\cap Q$ and $Q\supseteq P\cap Q$ (the union of hulls is monotone), so both contain every component with a contact on $C_i$, and so does $P\cap Q$. There are finitely many components, so the common refinement of all admissible partitions is admissible and is the finest. If a block $P$ of the finest partition is long on neither square and contains two overlap classes, replace $P$ by its overlap classes: (A1) still holds, the new blocks are long on neither square (they are subsets of $P$), and the other blocks are unchanged, so the new partition is admissible and strictly finer, a contradiction. If no component has a contact on $C_{3-i}$, condition (A2) on $C_{3-i}$ is void; if no overlap class is long the overlap classes form an admissible partition, hence the finest, and these are the classes of Section~\ref{sec:joint-extraction} for $C_i$ with the intervals $\mathcal H(K)$ there, which equal the unions of the members' hulls; if some overlap class is long, every admissible partition has a long block containing it, which by (A2) contains every component (each has a contact on $C_i$), so the finest partition is the single class of all components, the whole-contour group of Section~\ref{sec:joint-extraction}. The term of a group: let $\mathfrak P$ be the finest admissible partition of a term $T$ and $W\in\mathfrak P$. Every cluster occurrence of $W$ hits a root of $T$, which is joined to it by an edge and so lies in its component, hence in $W$; so $T_W$ is a term of \eqref{eq:c-cluster-expansion}. Edges are defined pairwise, so the graph of $T_W$ is the subgraph of that of $T$ induced on $W$, and $W$ being a union of components of $T$, the components of $T_W$ are those of $T$ in $W$, with the same contact sets, hulls and long-ness; by (A1) $W$ is a union of overlap classes of $T$, which are the overlap classes of $T_W$. Let $\mathfrak P'$ be an admissible partition of $T_W$ and $\mathfrak P''=\mathfrak P'\cup(\mathfrak P\setminus\{W\})$. (A1) holds for $\mathfrak P''$, and (A2) for its blocks outside $W$; if $Q\in\mathfrak P'$ is long on $C_i$, then $W\supseteq Q$ is long on $C_i$, so $W$ contains every component of $T$ with a contact on $C_i$ by (A2) for $\mathfrak P$, and all of them lie in $Q$ by (A2) for $\mathfrak P'$. So $\mathfrak P''$ is admissible for $T$, $\mathfrak P$ refines it, and $\mathfrak P'=\{W\}$. Long-ness and hulls being the same in $T$ and $T_W$, the interval set is the same.

(b) Distinct groups of a term have disjoint intervals: if neither is long on $C_i$, their $C_i$-parts are unions of hulls of components in different blocks, hence of different overlap classes by (A1), and such hulls are disjoint; if one is long on $C_i$, it contains every component with a contact on $C_i$, so the other has $\mathcal H_i=\varnothing$. Domains in different groups are compatible as in the proof of Lemma~\ref{lem:grouped-source} (no edge of the term's graph joins different components). Every component's hull satisfies $H^i_\gamma\subseteq\mathcal H_i(W)$ for its group $W$: trivially if $W$ is long on $C_i$, and by definition otherwise. Reversibility: let $\mathfrak B=\{W_1,\ldots,W_m\}$ be a compatible family, each $W_j$ being the only group of its own term, and let $\mathfrak P$ be the finest admissible partition of the union term, whose components are those of the $W_j$ (compatibility excludes edges between groups). $\mathfrak B$ is admissible for the union term: components of different $W_j$ have hulls in disjoint intervals, so they do not overlap and (A1) holds; if $W_j$ is long on $C_i$, then $\mathcal H_i(W_j)=\mathcal I_{C_i}$, so every other $W_k$ has $\mathcal H_i(W_k)=\varnothing$, no contact on $C_i$, and $W_j$ contains every component with a contact on $C_i$. Hence $\mathfrak P$ refines $\mathfrak B$. The restriction of $\mathfrak P$ to the components of $W_j$ is admissible for the term of $W_j$ alone: its overlap classes are those of the union term inside $W_j$, and a block of $\mathfrak P$ inside $W_j$ that is long on $C_i$ contains every component of the union term with a contact on $C_i$, in particular those of $W_j$. Since $\{W_j\}$ is the finest admissible partition of that term, $\mathfrak P$ restricted to $W_j$ is the single block $W_j$, so $\mathfrak P=\mathfrak B$, with the same intervals. The rest of the proof of Lemma~\ref{lem:grouped-source} (interleaving count, multiplicities of identical occurrences, absolute convergence) does not involve the intervals and applies verbatim.

(c) Let $W$ be a $C_i$-group of the pair. Its roots are marked and $\mathcal H_{3-i}(W)=\varnothing$, so every root $D$ has $\mathsf T_D=\mathsf T^i_D\neq\varnothing$, and the components, their hulls on $C_i$, the overlap graph and the long-ness on $C_i$ of the term of $W$ are the same for the pair and for the square $C_i$; admissibility for the pair reduces to admissibility for $C_i$ (condition (A2) on $C_{3-i}$ is void), so $W$ is the only group of its term for $C_i$ as well, with the same interval. Conversely a single-square group of $C_i$ whose roots have no contact on $C_{3-i}$ has the same term data for the pair. The weight \eqref{eq:part1-group-weight} and the footprint of Section~\ref{sec:c4-6} depend on the roots, the cluster list and the interval only. Compatibility for the pair of two $C_i$-groups is compatibility of domains and disjointness of the $C_i$-intervals, which is Definition~\ref{def:groups} for $C_i$ (its singleton clause for a whole-contour group being implied by the disjointness, since every group has a nonempty interval); for a $C_1$-group and a $C_2$-group the interval condition is void. (d) $W$ is not long on $C_{3-i}$, since $\mathcal H_{3-i}(W)=\varnothing$, so by (a) it is an overlap class; its members have hulls on $C_i$ only, so they are chained by intersections of these hulls, and $\mathcal H_i(W)$ is a union of chained cyclic intervals with fewer than $n_{C_i}/2$ positions in total, which is a cyclic interval by the argument of Section~\ref{sec:joint-extraction}.
\end{proof}

\subsection{Reference scalars and the joint extraction}\label{sec:two-loops-extraction}
For a $C_i$-group $W$ let $P^\circ_W(a)$ be the ordered product of the reference links along the consecutive traversals of $\mathcal H_i(W)$, in the cyclic order of $C_i$ starting from the first traversal of the interval (from the base point of $C_i$ if $W$ is long on $C_i$), as in Section~\ref{sec:joint-extraction}, and define $a_W$ by \eqref{eq:part1-reference-group} with $\chi_i$ in place of $\chi$; $\mathcal H_i(W)$ is a cyclic interval or $\mathcal I_{C_i}$ by Lemma~\ref{lem:two-loop-grouping}(d), and $a_W=a^{(i)}_W$ by Lemma~\ref{lem:two-loop-grouping}(c), since the reference links and the measures $\widehat\nu^{\mathrm f,\circ}_{\xi,D}$ do not involve the observable (Lemma~\ref{lem:probe-measures}). For a cross group put $a_W:=0$. With these numbers define, for the pair,
\begin{equation}\label{eq:two-loop-reference-gas}
    S^\circ_{\xi,C}=\sum_{\mathfrak B\text{ compatible}}\prod_{W\in\mathfrak B}a_W\;,
\end{equation}
in which a family containing a cross group contributes zero, the restricted sums $S^{\circ,[\mathfrak B]}_{\xi,C}$ over the families compatible with every member of $\mathfrak B$, and the joint defect $R_\xi(\mathfrak B)$ by \eqref{eq:part1-joint-defect}.

\begin{lemma}[Schur identity and joint extraction for two squares]\label{lem:two-loop-schur}
(a) For a $C_i$-group $W$ and $h\in G$, \eqref{eq:part1-closed-schur} holds with $\chi_i$ in place of $\chi$. (b) Proposition~\ref{prop:joint-extraction} holds for the pair: \eqref{eq:part1-scalar-extraction} and \eqref{eq:c-one-step-branch-identity} hold with $S^\circ_{\xi,C}$ of \eqref{eq:two-loop-reference-gas}, the sums converging absolutely. (c) At the reference data, $R^\circ_\xi(\mathfrak B)=0$ for every nonempty compatible family $\mathfrak B$ without cross groups.
\end{lemma}

\begin{proof}
(a) The proof of Lemma~\ref{lem:schur} for the square $C_i$: it uses only the invariance of the reference measures under the constant gauge rotations $R_u$, the covariance $P^\circ_W(R_ua)=uP^\circ_W(a)u^{-1}$ and Schur's lemma for $\pi_i$. (b) The proof of Proposition~\ref{prop:joint-extraction} is the Boolean inversion of \eqref{eq:part1-joint-defect} and an absolutely convergent rearrangement; it uses only $|a_W|\leq\beta_W$, which holds for $a_W=0$, the finiteness of $\sum_W\beta_W$ (a group is determined by its root family, cluster list and interval, Lemma~\ref{lem:vacuum-clusters} and Lemma~\ref{lem:c2-reference}) and Lemma~\ref{lem:two-loop-grouping}(b) in place of Lemma~\ref{lem:grouped-source}. (c) Let $\mathcal S$ be a compatible family of $C_1$- and $C_2$-groups. The contacts on $C_i$ of the roots of a $C_i$-group $W$ lie in $\mathcal H_i(W)$ (Lemma~\ref{lem:two-loop-grouping}(b)), and the roots of a $C_{3-i}$-group have no contact on $C_i$. So, as in the proof of Lemma~\ref{lem:schur}, a traversal outside every interval of $\mathcal S$ has the reference link $\id$, and a traversal in $\mathcal H_i(W)$ has at $\mathsf e_{\mathcal A(\mathcal S)}a$ its link at $\mathsf e_{\mathcal A(W)}a$, since it lies in no contact set of a root of another group; the links of $\mathcal H_i(W)$ multiply to $P^\circ_W(a)$. Number the $C_i$-groups of $\mathcal S$ in the cyclic order of $C_i$; by cyclicity of the trace the reference Wilson function at $\mathsf e_{\mathcal A(\mathcal S)}a$ is $\prod_{i=1,2}\chi_i\bigl(\prod_{W\in\mathcal S,\ C_i\text{-group}}P^\circ_W(a)\bigr)$, the product over the $C_i$-groups in that order (an empty product being $\id$), and a family containing a group long on $C_i$ has no other $C_i$-group. Integrate one group at a time by Fubini, with (a) for the group integrated, the product of the other $P^\circ$ of the same square as $h$, and the character of the other square as a constant factor: $A^\circ_\xi(\mathcal S)=\prod_{W\in\mathcal S}a_W$. Hence $R^\circ_\xi(\mathfrak B)=\prod_{W\in\mathfrak B}a_W\sum_{\mathcal S\subseteq\mathfrak B}(-1)^{|\mathfrak B\setminus\mathcal S|}=0$ for $\mathfrak B\neq\varnothing$.
\end{proof}

For a family containing a cross group, $R^\circ_\xi(\mathfrak B)$ need not vanish (for instance $R^\circ_\xi(\{W\})=A^\circ_\xi(\{W\})$ for a cross group $W$); the defects of such families are bounded directly (Lemmas~\ref{lem:two-loop-defect} and~\ref{lem:two-loop-cross}). In place of the second equality of \eqref{eq:part1-reference-gas}, Lemma~\ref{lem:two-loop-grouping}(b) at the reference data (its proof uses only the field-independent geometry, as in Section~\ref{sec:joint-extraction}) and the Fubini step of (c) give
\[
    \frac{\int f^\circ_\xi K^\circ_{\xi;0}\,\dd a}{I^\circ_{\xi;0}}=S^\circ_{\xi,C}+\sum_{\mathfrak B\text{ compatible, with a cross group}}A^\circ_\xi(\mathfrak B)\;.
\]
The choice $a_W:=0$ is made for every cross group, and no reduction of the reference integrals of cross families to products of numbers is used. The pair scalar is identified with the single-square scalars by Lemma~\ref{lem:two-loop-cross}(b).

\subsection{The defects of a family with cross groups}\label{sec:two-loops-defects}
For a $C_i$-group $W$ the closed contraction $c_W(U_1;h)$ of \eqref{eq:c4-6-2}, its defect $d_W$ and $\mathfrak d_W(U_1)$ of \eqref{eq:c4-6-3} are formed with $\chi_i$ and the interval $\mathcal H_i(W)$; by Lemma~\ref{lem:two-loop-grouping}(c) they are the single-square quantities of $C_i$. For every group, cross groups included, $M(W)$ is the positive majorant of Section~\ref{sec:c4-6-2}: the product of the root weights $\rho\,m(D)$, the cluster weights $M_v(\mathfrak c_j)$ and $1/k(W)!$, with $\rho$ of Lemma~\ref{lem:two-loop-source}; for $C_i$-groups this only enlarges the single-square majorant, so every single-square bound by $M(W)$ still holds. Put $K_\chi=\max(K_{\chi_1},K_{\chi_2})$ with the constants of Lemma~\ref{lem:c4-6-1}(a).

\begin{lemma}[sources with two closed words]\label{lem:two-loop-source}
Let $f(a)=\prod_{i=1,2}\dim(\pi_i)^{-1}\operatorname{tr}Y_i(a)$, where $Y_i(a)$ is an ordered product of the matrices $\pi_i(U_\xi(a)(r))^{\pm1}$ of a set of reconstructed contour links of $C_i$, each traversal used at most once, and of fixed unitary matrices of $V_{\pi_i}$. Then \eqref{eq:c4-4-3}, \eqref{eq:c4-10-8}--\eqref{eq:c4-10-9}, the undifferentiated bound \eqref{eq:c4-7-6} and the last statement of Lemma~\ref{lem:c4-7-2}, and Lemma~\ref{lem:c4-4-2} hold for $f$, with $C_\chi$ enlarged once to a constant of $(\chi_1,\chi_2)$ and $g_0$ so small that $u\leq u_*$ and $v\leq1$ with the enlarged constant; in particular $|\mathfrak Q_\xi(\mathcal A;f)|\leq\prod_{D\in\mathcal A}\rho\,m(D)$ for every nonempty compatible root family $\mathcal A$, with $\rho$ of \eqref{eq:c4-4-8} formed with the enlarged constant.
\end{lemma}

\begin{proof}
$f$ is the normalised trace of $Y_1(a)\otimes Y_2(a)$ on $V_{\pi_1}\otimes V_{\pi_2}$, and $|f|\leq1$; since $\operatorname{tr}(A\otimes B)=\operatorname{tr}A\operatorname{tr}B$, the normalised trace on $V_{\pi_1}\otimes V_{\pi_2}$ satisfies $|\dim^{-1}\operatorname{tr}(M_1\cdots M_k)|\leq\prod_j\|M_j\|_{\mathrm{op}}$. The proof of \eqref{eq:c4-10-8} uses only this inequality, the unitarity of the untouched factors, the row locality and the fixed derivative constants of the rows. For the row locality: $V(C_1)\cap V(C_2)=\varnothing$, so the bond sets of $C_1$ and $C_2$ are disjoint, each bond carries at most one traversal of $C_1\cup C_2$, and the number of reconstructed links a free label affects is the geometric count of bonds $r$ with the label in $\mathcal F_r$, at most $K$ as for one square, each traversal being used at most once in the two words together; the row input count, which fixes the derivative order of \eqref{eq:c4-10-4}, is unchanged. Then the product rule distributes the derivatives over the links of both words, with at most $K^q$ row assignments, and every untouched factor, a reconstructed link or a fixed matrix, is unitary in $V_{\pi_1}\otimes V_{\pi_2}$. The remainder bounds \eqref{eq:c4-10-9} follow by the same Taylor argument. The proofs of \eqref{eq:c4-7-6}, of the last statement of Lemma~\ref{lem:c4-7-2} and of Lemma~\ref{lem:c4-4-2} use, besides these bounds, only properties of the selected measures (the reversal of a root with empty ordinary selection, the Cauchy difference in $\zeta_j$, the majorants $M_\ell$ and \eqref{eq:c4-4-7}) and the vanishing of $\Delta_{\mathsf J}f$ when a selected label does not influence the words; none of them uses the form of $f$ otherwise.
\end{proof}

\begin{lemma}[defect factorisation for two squares]\label{lem:two-loop-defect}
On a branch covered for $C_1\sqcup C_2$ (Section~\ref{sec:two-loops-setting}), for every compatible family $\mathfrak B$ with cross groups $\mathfrak B_\times\subseteq\mathfrak B$,
\begin{equation}\label{eq:two-loop-defect}
    |R_\xi(\mathfrak B)|\leq\prod_{W\in\mathfrak B\setminus\mathfrak B_\times}K_\chi\mathfrak d_W(U_1)\prod_{W\in\mathfrak B_\times}M(W)\;.
\end{equation}
\end{lemma}

\begin{proof}
As in the proof of Lemma~\ref{lem:c4-6-1}(b), a traversal of $\mathcal H(W)$, $W\in\mathfrak B$, has no contact with a root of another group of $\mathfrak B$ (the intervals are disjoint and contain the contacts), so its reconstructed link depends only on the root variables of $W$, and a traversal outside every interval has its background link $U_1$ (Lemma~\ref{lem:covered}, Section~\ref{sec:c4-8-1}). For a $C_i$-group $W\in\mathfrak B$ let $X_W=w_\xi(W)\int\pi_i(T_WP_W^{-1})\,\dd\bigotimes_{D\in\mathcal A(W)}\widehat\nu^{\mathrm f}_{\xi,D}$ and $\Upsilon_W=X_W-a_W\id$, with $T_W$, $P_W$ of Section~\ref{sec:c4-6-1}. Since $a_W=0$ for cross groups, the term of $\mathcal S$ in \eqref{eq:part1-joint-defect} vanishes unless $\mathcal S\supseteq\mathfrak B_\times$. Fix the root variables $a_\times$ of the cross groups. For $i=1,2$ let $Y_i$ be the ordered product, in the cyclic order of $C_i$ starting at the first traversal of the interval of a $C_i$-group of $\mathfrak B$ if there is one and at the base point of $C_i$ otherwise (the trace is cyclic; as in the proof of Lemma~\ref{lem:c4-6-1}(b)), of the following factors: for a $C_i$-group $W\in\mathfrak B$, the matrix $\Upsilon_W\pi_i(P_W)$ at the position of $\mathcal H_i(W)$; for a traversal in $\mathcal H_i(W)$ of a cross group $W$, its reconstructed link at $a_\times$; for every other traversal, its background link. Expanding each $\Upsilon_W=X_W-a_W\id$ and using Fubini over the product measures of the $C_1$- and $C_2$-groups, exactly as in the proof of Lemma~\ref{lem:c4-6-1}(b), the term of $\mathcal S\supseteq\mathfrak B_\times$ is reproduced, because $f_\xi$ at $\mathsf e_{\mathcal A(\mathcal S)}a$ is the product of the two normalised traces and each $X_W$ enters exactly one of them; hence
\begin{equation*}
    R_\xi(\mathfrak B)=\Bigl(\prod_{W\in\mathfrak B_\times}w_\xi(W)\Bigr)\int\prod_{i=1,2}\frac{\operatorname{tr}Y_i(a_\times)}{\dim\pi_i}\,\dd\bigotimes_{W\in\mathfrak B_\times}\bigotimes_{D\in\mathcal A(W)}\widehat\nu^{\mathrm f}_{\xi,D}(a_\times)\;.
\end{equation*}
The integrand is multilinear in the matrices $\Upsilon_W$, $W\in\mathfrak B\setminus\mathfrak B_\times$. Write $\Upsilon_W=\|\Upsilon_W\|_{\mathrm{op}}T_W'$ with $\|T'_W\|_{\mathrm{op}}\leq1$, and each contraction $T'=V|T'|$ (polar decomposition, $V$ unitary) as $T'=\frac12(Ve^{iA}+Ve^{-iA})$ with $A=\arccos|T'|$, an average of two unitary matrices. The integrand is then $\prod_W\|\Upsilon_W\|_{\mathrm{op}}$ times an average of $2^{|\mathfrak B\setminus\mathfrak B_\times|}$ functions of the form of Lemma~\ref{lem:two-loop-source}, depending on $a_\times$ only through the links of the cross groups' traversals. If $\mathfrak B_\times\neq\varnothing$, that lemma applied to the compatible root family $\mathcal A(\mathfrak B_\times)$ bounds each such integral by $\prod_{D\in\mathcal A(\mathfrak B_\times)}\rho\,m(D)$ in modulus; if $\mathfrak B_\times=\varnothing$, there is no integral, and each of the unitary words has normalised traces of modulus at most one (for $\mathfrak B=\varnothing$ the bound is $|f_\xi(0)|\leq1$). Further $|w_\xi(W)|\leq\prod_jM_v(\mathfrak c_j)/k(W)!$: by \eqref{eq:part1-vacuum-cluster} the cluster weight is $\varphi^T/m!$ times the product of the activities, $|\varphi^T|$ is at most the number of spanning trees of the incompatibility graph (proof of Lemma~\ref{lem:vacuum-clusters}, Section~\ref{app:kp}), and the activities are bounded by $m(D)$ by \eqref{eq:fixed-kernel-raw-activity-payment} of Lemma~\ref{lem:activity-bounds}. Together with $\|\Upsilon_W\|_{\mathrm{op}}\leq K_{\chi_i}\mathfrak d_W(U_1)\leq K_\chi\mathfrak d_W(U_1)$ for a $C_i$-group (Lemma~\ref{lem:c4-6-1}(a) for the square $C_i$) this is \eqref{eq:two-loop-defect}; the bound holds with $K_{\chi_i}$ on the $C_i$-groups.
\end{proof}

\subsection{Pins for the groups of two squares}\label{sec:two-loops-pins}
For a group $W$ of the pair, $n(W)$, $\mathfrak s(W)=n(W)+|\mathcal H(W)|/(LM)$, the footprint $F(W)$, $X(W)$, $\pi_1(\mathcal H(W))$ and the envelope $\mathfrak E(W)$ are defined as in Section~\ref{sec:c4-6} and Lemma~\ref{lem:c4-6-env}, with $\mathcal H(W)=\mathcal H_1(W)\sqcup\mathcal H_2(W)$ and the traversals of both squares; for components, $t_\gamma$, $n_\gamma$, $X_\gamma$, $F(\gamma)$ and $\mathfrak s(\gamma)=n_\gamma+|H_\gamma|/(LM)$ are those of Section~\ref{sec:c4-8-2} with $H_\gamma=H^1_\gamma\sqcup H^2_\gamma$, and $\sigma_K=\sum_{\gamma\in K}\mathfrak s(\gamma)$. Put $A'_E=17(2R_{\mathrm e}+1)^4$.

\begin{lemma}[pins for two squares]\label{lem:two-loop-pins}
Under the hypotheses of Lemma~\ref{lem:c4-6-2}, with its constant $C$ enlarged by a fixed factor:
(a) Lemma~\ref{lem:c4-6-env} holds for every group $W$ of the pair with $16n(W)+|\mathcal H(W)|/(LM)$ and $17\mathfrak s(W)$ in place of the bounds $8+|\mathcal H(W)|/(LM)$ and $9\mathfrak s(W)$ of its part (a), $A'_E$ in place of $A_F$ and $A_E$, and $\operatorname{dist}(\Delta,C)$ the distance to $V(C_1)\cup V(C_2)$; if $W$ is a cross group, $\pi_1(\mathcal H(W))\subseteq\mathfrak E(W)$ contains a cube containing a site of $V(C_1)$ and a cube containing a site of $V(C_2)$.
(b) Lemma~\ref{lem:c4-8-2} holds for the components of the pair, with $|H^i_\gamma|\leq A_{\mathrm h}LMn_\gamma$ for $i=1,2$, the hull pin \eqref{eq:c4-8-5} for $a\in\mathcal I_{C_i}$ taken with the contour distance $d_{C_i}$ of $C_i$, and the sum over $a'\in\mathcal I_C$ in \eqref{eq:c4-8-6}.
(c) Lemma~\ref{lem:c4-8-3}(a) holds for the overlap classes of the pair. A class long on $C_i$ satisfies $\sigma_K\geq n_{C_i}/(2LM)$ and $\sum_{K\text{ long on }C_i}t_Ke^{34\eta\sigma_K}\leq C\rho n_{C_i}e^{-15\eta n_{C_i}/(LM)}$. The groups long on some square have total weighted mass $\sum_{W\text{ long}}M(W)e^{4\eta\mathfrak s(W)}\leq2C\rho\bigl(n_{C_1}e^{-10\eta n_{C_1}/(LM)}+n_{C_2}e^{-10\eta n_{C_2}/(LM)}\bigr)\leq C'\rho$, uniformly in $n_1,n_2$.
(d) The pins \eqref{eq:c4-6-5} hold for the groups of the pair, with $\mathcal I_C$ the traversals of both squares; for every group $W$ of the pair, cross groups included, $\sum_{W'\not\sim W}|a_{W'}|e^{3\eta\mathfrak s(W')}\leq\frac18\eta\mathfrak s(W)$, the sum over the groups of the pair; and Lemma~\ref{lem:c4-6-3} and Lemma~\ref{lem:c4-9-2} hold for the reference gas \eqref{eq:two-loop-reference-gas}, with $C\rho n_C$ in \eqref{eq:c4-9-4}.
\end{lemma}

\begin{proof}
(a) Count. On each square, $\mathcal H_i(W)$ is either $\mathcal I_{C_i}$ or, if $W$ is not long on $C_i$, a union of member hulls with fewer than $n_{C_i}/2$ positions in total; in the second case each member hull is a cyclic interval (none is the whole of $\mathcal I_{C_i}$), each maximal cyclic interval of the union has fewer than $n_{C_i}/2$ positions because the union does, and there are at most $|W|$ of them, where $|W|$, the number of components of $W$, is at most $n(W)=\sum_\gamma n_\gamma$ since $n_\gamma\geq1$. As in the proof of Lemma~\ref{lem:c4-6-env}(a), a cyclic interval of fewer than $n_{C_i}/2$ traversals meets each side in one run and lies in at most $8+t/(LM)$ cubes for $t$ traversals, and so does $\mathcal I_{C_i}$; summing over at most $2n(W)$ intervals, $|\pi_1(\mathcal H(W))|\leq16n(W)+|\mathcal H(W)|/(LM)$, so $|X(W)\cup\pi_1(\mathcal H(W))|\leq17\mathfrak s(W)$ and $|F(W)|\leq|\mathfrak E(W)|\leq A'_E\mathfrak s(W)$. Connectivity. For a component $\gamma$, $X_\gamma$ is connected (proof of Lemma~\ref{lem:c4-6-env}(b)); $\pi_1(H^i_\gamma)$ is connected, $H^i_\gamma$ being a cyclic interval or $\mathcal I_{C_i}$, and it contains the cube of the initial vertex of a contact of $\gamma$ on $C_i$, which is at index distance at most one from a cube of a root of $\gamma$; so $X_\gamma\cup\pi_1(H_\gamma)$ is connected. Two overlapping components share a traversal, hence the cube of its initial vertex. So $X(W)\cup\pi_1(\mathcal H(W))$ is connected if $W$ is an overlap class, which by Lemma~\ref{lem:two-loop-grouping}(a) covers every group long on neither square. Let $W$ be long on $C_i$. Each overlap class in $W$ has a component with a contact on $C_i$, or $W$ is also long on $C_{3-i}$ and the class has a contact on $C_{3-i}$: otherwise splitting this class off $W$ gives an admissible partition (the remainder of $W$ keeps every component with a contact on a square on which $W$ is long and keeps the same union of hulls there, the class is long on no square, being a subset of $W$ without contacts on those squares and with the union of its hulls on the other square inside that of $W$), contradicting the finest choice. If $W$ is long on $C_i$ only, every class in $W$ meets the connected set $\pi_1(\mathcal I_{C_i})$, so the union is connected. If $W$ is long on both squares and no class in $W$ has contacts on both, then splitting $W$ into its classes with contacts on $C_1$ and those with contacts on $C_2$ is admissible (each part contains every component with a contact on the square on which it is long, and has no contact on the other square), a contradiction; so some class of $W$ meets both $\pi_1(\mathcal I_{C_1})$ and $\pi_1(\mathcal I_{C_2})$, and the union is connected. Hence $\mathfrak E(W)$ is connected. The statements (c) and (d) of Lemma~\ref{lem:c4-6-env} then follow by its proof. For a cross group both $\mathcal H_1(W)$ and $\mathcal H_2(W)$ are nonempty, and $\pi_1(\mathcal H(W))$ contains the cubes of their vertices.

(b) The proof of Lemma~\ref{lem:c4-8-2}. Part (a) there does not involve hulls. In (b) the contacts of $\gamma$ on $C_i$ lie within $2L$ of $X_\gamma$, and the square geometry of $C_i$ (Section~\ref{sec:extent-weights}) gives $|H^i_\gamma|\leq A_{\mathrm h}LMn_\gamma$ and $d_{C_i}(a,a')\leq A_{\mathrm h}LMn_\gamma$ for $a\in H^i_\gamma$ and a contact $a'$ of $\gamma$ on $C_i$; the contour sum is over $\mathcal I_{C_i}$. In \eqref{eq:c4-8-6} the sum over the traversals of each square is bounded as there, and there are two squares. For part (c) of Lemma~\ref{lem:c4-8-2}: for the pair $\mathfrak s(\gamma)\leq(1+2A_{\mathrm h})n_\gamma\leq2A_{\mathrm s}n_\gamma$, so $100\eta\mathfrak s(\gamma)\leq200\eta A_{\mathrm s}n_\gamma<\varpi n_\gamma/40$ by $1000\eta A_{\mathrm s}<\varpi/8$, and the pins with the weight $e^{100\eta\mathfrak s(\gamma)}$ are dominated by those of (b) with the weight $e^{\varpi n_\gamma/4}$, with a margin $e^{\varpi n_\gamma/8}$; $\eta$ is not changed. In the footprint pin, at most $CLM$ traversals of each square have stencils meeting a given cube, so at most $2CLM$ in all, and the footprint pin holds with the constant doubled.

(c) The proof of Lemma~\ref{lem:c4-8-3}(a) applies with $H_\gamma=H^1_\gamma\sqcup H^2_\gamma$: the components whose hulls meet $H_\gamma$ have weighted mass at most $C\rho|H_\gamma|$ by the hull pins of (b) at the traversals of $H_\gamma$, and an overlap class long on neither square has the interval $\bigcup_\gamma H_\gamma$, of size at most $\sum_\gamma|H_\gamma|\leq LM\sigma_K$. A class long on $C_i$ has $\sigma_K\geq\sum_\gamma|H^i_\gamma|/(LM)\geq n_{C_i}/(2LM)$; anchoring it at a traversal of $\bigcup_\gamma H^i_\gamma$ ($n_{C_i}$ choices) and using the class pin with the weight $e^{64\eta\sigma_K}$ gives the stated sum, the exponent $30\eta\sigma_K$ being at least $15\eta n_{C_i}/(LM)$. Long groups: if no overlap class of a term is long, the overlap classes are admissible and no group is long. Let $W$ be long; it contains an overlap class $K_0$ long on some square $C_{i_0}$: otherwise replacing $W$ by its overlap classes keeps (A1), the new blocks are long on no square, and the other blocks are unchanged, so the new partition is admissible, and it is strictly finer unless $W$ is a single class, which is then not long; both contradict the choice of $W$. $W$ is long on $C_{i_0}$, and by the splitting argument of (a) it is the union of the overlap classes with a contact on a square on which it is long. As in Lemma~\ref{lem:c4-8-3}(c), the sum of $M(W)$ over the groups with a given set of components is at most $\prod_\gamma t_\gamma$, and $\mathfrak s(W)\leq\sum_{\gamma\in W}\mathfrak s(\gamma)+\sum_{i:\,W\text{ long on }C_i}n_{C_i}/(LM)$. If $W$ is long on $C_{i_0}$ only, the classes other than $K_0$ are anchored at traversals of $C_{i_0}$ and sum to at most $e^{C\rho n_{C_{i_0}}}$ with their weights $e^{64\eta\sigma}$, the interval weight is $e^{4\eta n_{C_{i_0}}/(LM)}$, and $K_0$ contributes $C\rho n_{C_{i_0}}e^{-15\eta n_{C_{i_0}}/(LM)}$; with $C\rho LM\leq\eta$ the product is at most $C\rho n_{C_{i_0}}e^{-10\eta n_{C_{i_0}}/(LM)}$. If $W$ is long on both squares, let $j\neq i_0$; since $\sum_{\gamma\in W}|H^j_\gamma|\geq n_{C_j}/2$ and $|H^j_\gamma|\leq LM\mathfrak s(\gamma)$, $1\leq e^{-15\eta n_{C_j}/(LM)}\prod_{\gamma\in W}e^{30\eta\mathfrak s(\gamma)}$; inserting this factor, the classes other than $K_0$ are anchored at traversals of $C_1$ or $C_2$ and sum to at most $e^{C\rho n_C}$ with the weights $e^{64\eta\sigma}$, the interval weight is $e^{4\eta n_C/(LM)}$, and $K_0$ contributes $C\rho n_{C_{i_0}}e^{-15\eta n_{C_{i_0}}/(LM)}$ with the weight $e^{34\eta\sigma_{K_0}}$; the product is at most $C\rho n_{C_{i_0}}e^{-10\eta n_C/(LM)}\leq C\rho n_{C_{i_0}}e^{-10\eta n_{C_{i_0}}/(LM)}$. The groups long on $C_i$ only thus contribute at most $C\rho n_{C_i}e^{-10\eta n_{C_i}/(LM)}$, and the groups long on both squares at most that for each $i_0=i$; summing over the four pairs (set of long squares, $i_0$) gives the bound, and $xe^{-10\eta x/(LM)}\leq LM/(10e\eta)$ for $x=n_{C_1},n_{C_2}$ gives $C'=2CLM/(5e\eta)$.

(d) The pins over groups long on neither square are the class pins of (c), as in Lemma~\ref{lem:c4-8-3}: such a group is an overlap class $K$ (Lemma~\ref{lem:two-loop-grouping}(a)), the sum of $M(W)$ over the groups with class $K$ is at most $t_K$, $\mathfrak s(W)\leq\sigma_K$ and $F(W)\subseteq\bigcup_{\gamma\in K}F(\gamma)$; the long groups are bounded by their total mass (c). For the second statement, a group $W'\not\sim W$ has a footprint cube in $\mathfrak N_1(F(W))$ or a traversal in $\mathcal H(W)$ by (a) and Lemma~\ref{lem:c4-6-env}(d), so the pins give at most $C\rho(3^4|F(W)|+|\mathcal H(W)|)\leq C\rho(3^4A'_E+LM)\mathfrak s(W)$, which is at most $\eta\mathfrak s(W)/8$ under $C\rho LM\leq\eta/8$ with the enlarged constant (and $e^{3\eta\mathfrak s}\leq e^{4\eta\mathfrak s}$, $|a_{W'}|\leq M(W')$). The proofs of Lemmas~\ref{lem:c4-6-3} and~\ref{lem:c4-9-2} then apply verbatim; the species of the reference gas are the groups with $a_W\neq0$, all of them $C_1$- or $C_2$-groups, and the exclusion ratio \eqref{eq:c4-6-7} holds for every compatible family $\mathfrak B$ of groups of the pair, cross groups included, because the criterion was proved for every $W$.
\end{proof}

\subsection{The cross terms}\label{sec:two-loops-cross}
Covered and far histories are those of Section~\ref{sec:two-loops-setting}. A single-square group $W$ of $C_i$ touches $C_{3-i}$ if some root $D$ of $W$ has $\mathsf T_D\cap\mathcal I_{C_{3-i}}\neq\varnothing$, with the contact set $\mathsf T_D$ of the pair. Put $c_\times=2\eta/A'_E$.

\begin{lemma}[the cross terms]\label{lem:two-loop-cross}
Under the hypotheses of Lemma~\ref{lem:two-loop-pins} and, for (c), those of Lemma~\ref{lem:c4-6-4} for each square, with a constant $C_\times$ fixed before the coupling and independent of $n_1$, $n_2$ and the positions of the squares, and with $C_\chi=\max_iC_{\chi_i}$ and $c=\min_ic_i$ formed from the constants $C_{\chi_i}$, $c_i$ of \eqref{eq:c4-6-8} for the square $C_i$:
(a) For every cross group $W$ of the pair and every $a\in\mathcal H(W)$, $\mathfrak s(W)>\delta_a/(A'_ELM)$; and $\sum_{W\text{ cross}}M(W)e^{\mathfrak a_W/8}\leq C_\times\rho\Delta_\times$.
(b) Let $\mathscr L$ and $\mathscr L_i$ be the connected-cluster logarithms of the reference gas \eqref{eq:two-loop-reference-gas} of the pair and of the reference gas of the single square $C_i$ (Lemma~\ref{lem:c4-9-2}), and $E^\times_\xi=\mathscr L-\mathscr L_1-\mathscr L_2$. Then
\begin{equation}\label{eq:two-loop-factorisation}
    S^\circ_{\xi,C}=S^\circ_{\xi,C_1}S^\circ_{\xi,C_2}e^{E^\times_\xi}\;,\qquad|E^\times_\xi|\leq C_\times\rho\Delta_\times\;.
\end{equation}
(c) On covered far histories, with $\mathscr R_\xi$ the family sum \eqref{eq:c4-1-1} of the pair,
\begin{equation}\label{eq:two-loop-family-sum}
    |\mathscr R_\xi|\leq|S^\circ_{\xi,C}|\Bigl[\exp\Bigl(C_\chi n_C\rho\bigl(\varepsilon_1+e^{-cb/(LM)}\bigr)+C_\times\rho\Delta_\times\Bigr)-1\Bigr]\;,\qquad|S^\circ_{\xi,C}|\leq e^{C_\times\rho\Delta_\times}\;.
\end{equation}
\end{lemma}

\begin{proof}
Distance from an envelope. Let $\mathfrak X$ be a finite set of cubes, connected under adjacency, with $|\mathfrak X|\leq k$. As in the proof of Lemma~\ref{lem:c4-6-env}(c), any two sites in cubes of $\mathfrak X$ are at periodic fine $\ell^\infty$ distance less than $kLM$. Hence if $\mathfrak X$ contains the cube of the initial vertex of a traversal $a\in\mathcal I_{C_i}$ and a cube containing a site of $V(C_{3-i})$, then $\delta_a<kLM$.

(a) For a cross group $W$, $\mathfrak E(W)$ is connected, has at most $A'_E\mathfrak s(W)$ cubes and contains the cubes of the initial vertices of the traversals of $\mathcal H(W)$ and a cube containing a site of $V(C_1)$ and one containing a site of $V(C_2)$ (Lemma~\ref{lem:two-loop-pins}(a)); so $\delta_a<A'_ELM\mathfrak s(W)$ for every $a\in\mathcal H(W)$. Every group has a nonempty interval, so, pinning a cross group at a traversal of its interval and using $\mathfrak a_W/8\leq\eta\mathfrak s(W)$, $e^{-3\eta\mathfrak s(W)}\leq e^{-3\eta\delta_a/(A'_ELM)}\leq e^{-c_\times\delta_a/(LM)}$ and the pin \eqref{eq:c4-6-5} of the pair (Lemma~\ref{lem:two-loop-pins}(d)),
\[
    \sum_{W\text{ cross}}M(W)e^{\mathfrak a_W/8}\leq\sum_{a\in\mathcal I_C}e^{-c_\times\delta_a/(LM)}\sum_{W\,|\,a\in\mathcal H(W)}M(W)e^{4\eta\mathfrak s(W)}\leq C\rho\Delta_\times\;.
\]

(b) The three gases. The species of the pair gas with nonzero activity are its $C_1$- and $C_2$-groups; the species of the gas of $C_i$ are the single-square groups of $C_i$, with the activities $a^{(i)}_W$. By Lemma~\ref{lem:two-loop-grouping}(c) and Lemma~\ref{lem:two-loop-schur}, the $C_i$-groups of the pair are the single-square groups of $C_i$ that do not touch $C_{3-i}$, with the same activities, and the compatibility relation of the pair restricted to them is that of the square $C_i$. Each gas satisfies the Kotecký--Preiss criterion in the following weighted form: for every species $W$ of that gas,
\begin{equation}\label{eq:two-loop-weighted-kp}
    \sum_{W'\not\sim W}|a_{W'}|e^{3\eta\mathfrak s(W')}\leq\frac18\eta\mathfrak s(W)\;,
\end{equation}
for the pair by Lemma~\ref{lem:two-loop-pins}(d), and for the square $C_i$ by the proof of \eqref{eq:c4-6-6}, which bounds $e^{\mathfrak a_{W'}}$ by $e^{4\eta\mathfrak s(W')}$ and so bounds $e^{3\eta\mathfrak s(W')}$ in the same way. This is the criterion of Section~\ref{app:kp} for the weights $\widetilde z_W=|a_W|e^{2\eta\mathfrak s(W)}$ and sizes $\eta\mathfrak s(W)$, so the rooted-tree argument there gives, for every species $W$ of each gas,
\begin{equation}\label{eq:two-loop-pinned-cluster}
    \sum_{\mathfrak c}m_W(\mathfrak c)|\Psi(\mathfrak c)|e^{2\eta\sum_{W'\in\mathfrak c}\mathfrak s(W')}\leq|a_W|e^{3\eta\mathfrak s(W)}\;,
\end{equation}
the sum over the cluster tuples $\mathfrak c$ of that gas, $\Psi(\mathfrak c)=\frac{\varphi^T(\mathfrak c)}{m!}\prod_{W'\in\mathfrak c}a_{W'}$, $m_W(\mathfrak c)$ the multiplicity of $W$ in $\mathfrak c$. The logarithms $\mathscr L$, $\mathscr L_i$ are the absolutely convergent series $\sum_{\mathfrak c}\Psi(\mathfrak c)$ of the three gases, $S^\circ_{\xi,C}=e^{\mathscr L}$ and $S^\circ_{\xi,C_i}=e^{\mathscr L_i}$ (Lemma~\ref{lem:c4-9-2}, for the pair by Lemma~\ref{lem:two-loop-pins}(d)), which gives the identity in \eqref{eq:two-loop-factorisation}.

Cancellation. A cluster tuple of the pair gas all of whose members are $C_i$-groups is a cluster tuple of the gas of $C_i$ with the same $\Psi$: the Ursell coefficient depends only on the incompatibility graph of the tuple, which is the same for both gases, and the activities agree. Conversely a tuple of the gas of $C_i$ none of whose members touches $C_{3-i}$ is such a tuple of the pair gas. Hence $E^\times_\xi=\sum_{\mathfrak c\in\mathcal X_0}\Psi(\mathfrak c)-\sum_{i=1,2}\sum_{\mathfrak c\in\mathcal X_i}\Psi(\mathfrak c)$, where $\mathcal X_0$ is the set of tuples of the pair gas containing a $C_1$-group and a $C_2$-group, and $\mathcal X_i$ the set of tuples of the gas of $C_i$ containing a group that touches $C_{3-i}$; the rearrangement is justified by absolute convergence.

Extent of the bridging tuples. Let $\mathfrak c\in\mathcal X_0\cup\mathcal X_1\cup\mathcal X_2$, let $W\in\mathfrak c$ and $a\in\mathcal H(W)$, $a\in\mathcal I_{C_i}$. The union $\mathfrak X$ of the envelopes of the members of $\mathfrak c$ is connected: each envelope is connected (Lemma~\ref{lem:two-loop-pins}(a) for the pair, Lemma~\ref{lem:c4-6-env}(b) for a single square), the incompatibility graph of a tuple with $\varphi^T(\mathfrak c)\neq0$ is connected, and two incompatible members have footprints at index distance at most one or a common traversal in their intervals (Lemma~\ref{lem:c4-6-env}(d), for the pair by Lemma~\ref{lem:two-loop-pins}(a)), the footprints and the cubes of the traversal vertices lying in the envelopes. $\mathfrak X$ contains the cube of the initial vertex of $a$, and a cube containing a site of $V(C_{3-i})$: for $\mathfrak c\in\mathcal X_0$, the cube of a traversal vertex of the interval of a member that is a $C_{3-i}$-group; for $\mathfrak c\in\mathcal X_j$ (then $i=j$), the cube of the initial vertex of a contact $n\in\mathsf T_D\cap\mathcal I_{C_{3-i}}$ of a root $D$ of a member $W_0$ touching $C_{3-i}$, which is within $2L$ of $D$ (proof of Lemma~\ref{lem:c4-6-env}: the stencil bond of the contact has its endpoints in $D$), hence at index distance at most one from a cube of $X(W_0)$ and in $\mathfrak E(W_0)$, since $R_{\mathrm e}\geq1$. With $|\mathfrak X|\leq A'_E\sum_{W'\in\mathfrak c}\mathfrak s(W')$ ($A_E\leq A'_E$), the distance bound gives $\sum_{W'\in\mathfrak c}\mathfrak s(W')>\delta_a/(A'_ELM)$.

Bound. Every member of a tuple has a nonempty interval. Pinning each bridging tuple at a member $W$ and a traversal $a\in\mathcal H(W)$, by \eqref{eq:two-loop-pinned-cluster} and the extent bound,
\[
    \sum_{\mathfrak c\in\mathcal X_0}|\Psi(\mathfrak c)|\leq\sum_{a\in\mathcal I_C}e^{-2\eta\delta_a/(A'_ELM)}\sum_{W\,|\,a\in\mathcal H(W)}\sum_{\mathfrak c}m_W(\mathfrak c)|\Psi(\mathfrak c)|e^{2\eta\sum_{W'\in\mathfrak c}\mathfrak s(W')}\leq\sum_{a\in\mathcal I_C}e^{-c_\times\delta_a/(LM)}\sum_{W\,|\,a\in\mathcal H(W)}M(W)e^{3\eta\mathfrak s(W)}\;,
\]
the inner sums over the species of the pair gas, and the pin \eqref{eq:c4-6-5} of the pair bounds the last sum by $C\rho\Delta_\times$. The same computation in the gas of $C_i$, with the single-square pin \eqref{eq:c4-6-5} for $C_i$ and the traversals of $\mathcal I_{C_i}$, bounds $\sum_{\mathfrak c\in\mathcal X_i}|\Psi(\mathfrak c)|$ by $C\rho\sum_{a\in\mathcal I_{C_i}}e^{-c_\times\delta_a/(LM)}$. Summing, $|E^\times_\xi|\leq2C\rho\Delta_\times$. The series are those of the countable index sets; their absolute convergence is Lemma~\ref{lem:c4-9-2}, and no truncation is needed.

(c) By Lemma~\ref{lem:two-loop-schur}(b) the family sum of the pair is \eqref{eq:c4-1-1}, $\mathscr R_\xi=\sum_{\mathfrak B\neq\varnothing}S^{\circ,[\mathfrak B]}_{\xi,C}R_\xi(\mathfrak B)$ over the compatible families of groups of the pair. By the ratio bound \eqref{eq:c4-6-7} for the pair (Lemma~\ref{lem:two-loop-pins}(d), valid for families containing cross groups) and Lemma~\ref{lem:two-loop-defect},
\[
    \bigl|S^{\circ,[\mathfrak B]}_{\xi,C}R_\xi(\mathfrak B)\bigr|\leq|S^\circ_{\xi,C}|\prod_{i=1,2}\ \prod_{W\in\mathfrak B,\ C_i\text{-group}}K_{\chi_i}\mathfrak d_W(U_1)e^{\mathfrak a_W/8}\prod_{W\in\mathfrak B,\ \text{cross}}M(W)e^{\mathfrak a_W/8}\;,
\]
with $K_{\chi_i}$ on the $C_i$-groups (the last sentence of the proof of Lemma~\ref{lem:two-loop-defect}). The $C_i$-groups of the pair are single-square groups of $C_i$ with the same $\mathfrak d_W$ and $\mathfrak a_W$, and the pair's $\rho$ dominates that of the square $C_i$ (\eqref{eq:c4-4-8} is increasing in $u$), so Lemma~\ref{lem:c4-6-4} for $C_i$ (on histories covered and far for both squares) bounds the sum of the first factors over the $C_i$-groups by $C_{\chi_i}n_{C_i}\rho(\varepsilon_1+e^{-c_ib/(LM)})$; the sum over all non-cross groups is at most $C_\chi n_C\rho(\varepsilon_1+e^{-cb/(LM)})$, and (a) bounds the sum of the second factors. Dropping the compatibility constraints, the sum over nonempty families is at most the exponential of the one-group sum minus one, as in the proof of Lemma~\ref{lem:c4-6-5}. Finally $|S^\circ_{\xi,C_i}|\leq1$ (Section~\ref{sec:c4-9-4}: a normalised integral of a normalised character against a positive kernel), so $|S^\circ_{\xi,C}|\leq e^{|E^\times_\xi|}$ by (b). The constant $C_\times$ is the largest of those above.
\end{proof}

\begin{corollary}[the scalar of the pair]\label{cor:two-loop-scalar}
Under the hypotheses of Lemma~\ref{lem:two-loop-cross} and of Proposition~\ref{prop:c4c-5-1} for each square, on $A_{\mathrm{cov}}$,
\begin{equation}\label{eq:two-loop-scalar}
    S^\circ_{\xi,C}=\kappa_{C_1}\kappa_{C_2}\,e^{E^\circ_{\xi,1}+E^\circ_{\xi,2}+E^\times_\xi}\;,
\end{equation}
with $E^\circ_{\xi,i}$ the error \eqref{eq:c4c-5-2} of the square $C_i$, bounded by \eqref{eq:c4c-5-1} with $n_i$, $J_{C_i}$ and the carrier distances $d_j$ of the segments of $C_i$, and $|E^\times_\xi|\leq C_\times\rho\Delta_\times$.
\end{corollary}

\begin{proof}
On $A_{\mathrm{cov}}$ the branch is covered for each square, so Proposition~\ref{prop:c4c-5-1} gives $S^\circ_{\xi,C_i}=\kappa_{C_i}e^{E^\circ_{\xi,i}}$; insert this in \eqref{eq:two-loop-factorisation}.
\end{proof}

\subsection{The retreated carrier}\label{sec:two-loops-retreated}
On a labelled term with final region $\Lambda$ (Section~\ref{sec:one-step-insertion}) the carrier is $D(\Lambda)$, a function of $\Lambda$ alone (Lemma~\ref{lem:c5-2-5}), so both squares share it. The groups of the pair at $D(\Lambda)$ are those of Definition~\ref{def:two-loop-groups} formed from the domains for $D(\Lambda)$ and the contact sets $\mathsf T_Z$ of Lemma~\ref{lem:c5-5-3}(a), taken over $n\in\mathcal I_C$; $a^\Lambda_W$ is the reference scalar of Lemma~\ref{lem:c5-5-3}(b) with $\chi_i$ for a $C_i$-group and $a^\Lambda_W:=0$ for a cross group, and $S^\circ_{\Lambda,C}:=\sum_{\mathfrak B\text{ compatible}}\prod_{W\in\mathfrak B}a^\Lambda_W$. $S^\circ_{\Lambda,C_i}$ is the labelled scalar of the square $C_i$, and $\mathscr R^\Lambda=\sum_{\varnothing\neq\mathfrak B\text{ compatible}}S^{\circ[\mathfrak B]}_{\Lambda,C}R^0_\Lambda(\mathfrak B)$ is the labelled defect \eqref{eq:c5-6-2} of the pair. $A^{\mathrm{fin}}_{\mathrm{cov}}$ (Section~\ref{sec:c5-3-4}) and $\mathsf N_b$ \eqref{eq:c4-2-1} are taken for $C=C_1\sqcup C_2$; the first is defined bond by bond and the second through $\Gamma_C=\Gamma_{C_1}\cup\Gamma_{C_2}$, so both are the intersections of the single-square events.

\begin{lemma}[the pair at the retreated carrier]\label{lem:two-loop-retreated}
Under the hypotheses of Proposition~\ref{prop:c5-6-def} and of \eqref{eq:c5-5-T} for each square, with the constants of Lemmas~\ref{lem:two-loop-source} and~\ref{lem:two-loop-pins}, on every labelled term, for exterior data in the support of the term's collar indicators and exterior cutoffs:
(a) $S^\circ_{\Lambda,C}=S^\circ_{\Lambda,C_1}S^\circ_{\Lambda,C_2}e^{E^\times_\Lambda}$ with $|E^\times_\Lambda|\leq C_\times\rho\Delta_\times$;
(b) on $A^{\mathrm{fin}}_{\mathrm{cov}}\cap\mathsf N_b^c$, $\mathscr R^\Lambda$ satisfies \eqref{eq:two-loop-family-sum} with $S^\circ_{\Lambda,C}$ in place of $S^\circ_{\xi,C}$;
(c) under \eqref{eq:c5-7-T} for each square, on $A^{\mathrm{fin}}_{\mathrm{cov}}$, $S^\circ_{\Lambda,C}=\kappa_{C_1}\kappa_{C_2}e^{E^\circ_{\Lambda,1}+E^\circ_{\Lambda,2}+E^\times_\Lambda}$, with $E^\circ_{\Lambda,i}$ the error of Corollary~\ref{cor:c5-7-2} for the square $C_i$.
\end{lemma}

\begin{proof}
(i) Lemma~\ref{lem:two-loop-grouping} and Lemma~\ref{lem:two-loop-pins}(a) do not involve the carrier: they use the contact sets, which are field- and $t$-independent (Lemma~\ref{lem:c5-5-3}(a)), and the footprints and hulls of the domains (Lemma~\ref{lem:c5-5-8}). (ii) Lemma~\ref{lem:two-loop-schur} at $D(\Lambda)$: the Schur identity at the reference data of Lemma~\ref{lem:c5-5-3}(b) for each $\pi_i$, as in the proof of Lemma~\ref{lem:two-loop-schur}(a), and the Boolean inversion with $a^\Lambda_W:=0$ for cross groups, Lemma~\ref{lem:c5-5-3}(c) with Lemma~\ref{lem:two-loop-grouping}(b) in place of Lemma~\ref{lem:grouped-source}. (iii) Lemma~\ref{lem:two-loop-source} at $D(\Lambda)$, at exterior parameter zero, on $A^{\mathrm{fin}}_{\mathrm{cov}}$: its proof uses $|f|\leq1$, the normalised-trace inequality, the row locality, the unitarity of the untouched factors and properties of the selected measures; the last are supplied at $D(\Lambda)$ by the paragraphs ``Charts and primitives'' and ``Parity and the source'' of the proof of Proposition~\ref{prop:c5-6-def} (even cutoffs $\chi_b$ and density \eqref{eq:fixed-kernel-density} at exterior parameter zero, the dummy comparison \eqref{eq:c4-4-7} as Lemma~\ref{lem:c5-5-2}(c), the majorants, and the source remainder bounds of Section~\ref{sec:c4-10-4} on $A^{\mathrm{fin}}_{\mathrm{cov}}$, where every contour row has its free inputs in $D(\Lambda)$ and no exterior variable enters a link, Lemma~\ref{lem:c5-5-7z}); those paragraphs use the source only through the properties listed in the proof of Lemma~\ref{lem:two-loop-source}. So $|\mathfrak Q^0_\Lambda(\mathcal A;f)|\leq\prod_{Z\in\mathcal A}\rho\,m(Z)$ for the two-word $f$. (iv) Lemma~\ref{lem:two-loop-defect} at $D(\Lambda)$: on $A^{\mathrm{fin}}_{\mathrm{cov}}$ the contour links are reconstructed from the young variable on $D(\Lambda)$ at the background $U_1$ (proof of Proposition~\ref{prop:c5-6-def}(a), Lemma~\ref{lem:c5-6-1}), so a traversal outside the intervals has the link $U_1$ and one in $\mathcal H(W)$ depends on the roots of $W$ only; the weights obey $|w^0_\Lambda(W)|\leq\prod_j|\Phi|^+_\Lambda(\mathfrak c_j)/k(W)!$ (Corollary~\ref{cor:c5-4-3}, Lemma~\ref{lem:c5-5-6}); with (iii) the proof gives \eqref{eq:two-loop-defect} with $\mathfrak d^0_W$ and $M^{(\rho)}_\Lambda(W)$. (v) Lemma~\ref{lem:two-loop-pins}(b)--(d) at $D(\Lambda)$: these are the summations of Lemmas~\ref{lem:c4-8-2} and~\ref{lem:c4-8-3} with the anchoring of the pair, and they use only the root weights $\rho\,m(Z)$, the tree weights $|\Phi|^+_\Lambda=M_v$ and the field-independent domain, cluster and interval geometry, which is the reason Lemma~\ref{lem:c5-5-8}(a) gives for one square; so they hold for $M^{(\rho)}_\Lambda$ with the constants of Lemma~\ref{lem:two-loop-pins}, and, since $\sigma(W)=\mathfrak s(W)$, $\nu=\eta$ and $C_F=A_E$ there, $c_\times$ is unchanged. In particular the weighted criterion \eqref{eq:two-loop-weighted-kp} holds at $D(\Lambda)$ for the pair, from $|a^\Lambda_W|\leq M^{(\rho)}_\Lambda(W)$ and the pins with $3^4A'_E+LM$ in place of $3^4A_F+LM$, and for each square as in the proof of Lemma~\ref{lem:two-loop-cross}(b). (vi) With (i)--(v), the proof of Lemma~\ref{lem:two-loop-cross}(a),(b) gives (a); for (b), the proof of Lemma~\ref{lem:two-loop-cross}(c) uses \eqref{eq:c5-6-def-3} for each square (Lemma~\ref{lem:c4-6-4} at $D(\Lambda)$ in the form with $e^{\mathfrak a_W/8}$, proved in the paragraph ``Pins'' of the proof of Proposition~\ref{prop:c5-6-def}) and $|S^\circ_{\Lambda,C_i}|\leq1$ (Lemma~\ref{lem:c5-6-3}). (c) is (a) with Corollary~\ref{cor:c5-7-2}(a) for each square, $S^\circ_{\Lambda,C_i}=\kappa_{C_i}e^{E^\circ_{\Lambda,i}}$ on $A^{\mathrm{fin}}_{\mathrm{cov}}$.
\end{proof}
